% =============================================================================
% Supplementary sections moved out of the main paper: E1 (controls),
% E2 (clinical proxy metrics), E3 (comparison to prior methods), E4 (lock-width sensitivity).
% =============================================================================

\section*{E. Additional Details}
\label{sec:suppl-details}

% -----------------------------------------------------------------------------
\subsection*{E1. Controls the comparison requires}
\label{sec:suppl-controls}

A control-arm comparison is only as trustworthy as the plumbing beneath it, and three preconditions we had assumed did not hold during development. (i) A single dataset object was shared by LoRA fine-tuning, the PPO pool and the metrics scan, so evaluation images had been seen during fine-tuning. We replaced it with one deterministic, persisted, \emph{patient-level} split shared by every stage, computed once because generation runs as independent array processes that cannot recompute an RNG-based split consistently. It assigns 80\% of patients to training and 20\% to test (seed 0; one image per patient in RSNA, so 21{,}347 training and 5{,}337 test images) and is stored with a fingerprint of the dataset. In every reported arm, generator fine-tuning, PPO rollouts and the validity probe use only the training fold, and every evaluated image is in the test fold. (ii) Training was unseeded: two identical submissions produced two different policies, which is survivable for a headline number but not for a comparison in which arms must differ by treatment alone. Every reported run now uses split seed 0, LoRA seed 0 and PPO seeds 1 (n$\to$p) and 2 (p$\to$n); evaluation uses base seed 0, with per-image seed equal to the base plus the image's global index ($+10^6$ for p$\to$n). Seeding fixes the random streams but not floating-point nondeterminism, so repeated runs are not guaranteed to be bit-identical. (iii) The control arms had never been exercised. Wiring them up surfaced a routing bug under which every reduce silently fell back to the trained arm's shards, which would have made ``the policy contributes nothing'' an artefact of a missing flag, and a shared output directory under which each arm overwrote the previous one's metrics.

A CPU-only invariant suite pins the LHR calibration, split disjointness and three further properties whose silent failure would yield a plausible-looking wrong number rather than a crash. Every training and evaluation run writes a manifest giving its code revision and configuration, and every evaluation shard carries a configuration stamp from which the transfer-ratio denominator is read (main paper Sec.~4.3). We report these because each was, at the time it was found, invisible in the outputs: all three would have produced a well-formed table.

% -----------------------------------------------------------------------------
\subsection*{E2. Clinical proxy metrics}
\label{sec:suppl-clinical}

We report two clinically motivated proxies under the committed configuration: the Cardio-Thoracic Ratio (CTR), the ratio of cardiac to thoracic width from the heart and lung segmentation masks, and a costophrenic-recess score (SCP) approximating pleural blunting, each tested with a paired $t$-test on the within-image delta (real minus counterfactual).

CTR shifts by $+0.0076$ ($t=2.246$, $p=0.026$) in n$\to$p and by $+0.0080$ ($t=1.010$, $p=0.314$) in p$\to$n: a small, only one-directionally significant shift, with an absolute magnitude ($<0.01$) far below the metric's clinical resolution (its diagnostic threshold is 0.5). This is consistent with the intended behaviour of the background lock, since the cardiac silhouette lies largely outside the lung mask that gates editing; we read it as a shift too small to act on either way, not as a clinically directional CTR effect. % TODO(J4)
SCP's own computation flags itself as degenerate under this configuration (relative spread 0.19\% across the four image groups), so its delta and $t$-test are withheld rather than reported on what is, numerically, a near-constant signal. Since it is built on the feature norm of a COCO-pretrained detector rather than a dedicated recess detector, we do not consider it informative in its current form. Neither proxy supports a directional claim about pleural or cardiac status. % TODO(J4)

% -----------------------------------------------------------------------------
\subsection*{E3. Comparison to prior methods}
\label{sec:suppl-comparison}

We do not compare head-to-head with prior counterfactual generators. None was evaluated under our split, seeding and evaluation protocol, and published values are not comparable to ours: they are computed with other classifiers, feature extractors and evaluation groups, and none reports LHR or held-out transfer. We therefore report no literature reference values alongside our measurements.

% -----------------------------------------------------------------------------
\subsection*{E4. Lock-width sensitivity (exploratory)}
\label{sec:suppl-lockwidth}

The width of the background lock's editable band (main paper Sec.~3.2: dilation and Gaussian feathering of the lung mask, both in latent cells) was set once and not tuned. As an exploratory check we regenerated the first 50 images of each direction with the committed policies under three settings: d0 (dilation 0, no feathering), d2 (dilation 2, feathering 1.0; the committed setting, so a rerun) and d4 (dilation 4, feathering 2.0). LHR uses the 47 annotated p$\to$n images. Table~\ref{tab:suppl-lockwidth} gives each setting's paired difference from the d2 rerun, with 95\% paired percentile-bootstrap CIs (10{,}000 resamples of the images, drawn jointly across settings; transfer as a ratio of means; CV also with an exact McNemar test). As the rerun-noise reference we use the difference between the d2 rerun and the committed run on the same images, and count a difference as exceeding rerun noise when its CI lies entirely outside $\pm$ that reference. The reference itself has CIs excluding zero on several metrics, so a CI that excludes zero is not on its own evidence of a setting effect. Every contrast here is across GPU types, and for p$\to$n metrics the committed reference was evaluated with the n$\to$p adapter, so its band also contains that change.

Narrowing the lock (d0) raises edit concentration ($+0.028$ n$\to$p, $+0.035$ p$\to$n) and lowers the fraction of pixels changed ($-0.060$, $-0.065$); it also moves two n$\to$p validator ratios beyond noise, in opposite directions (\texttt{nih} $-0.059$, \texttt{pc} $+0.019$). Widening it (d4) lowers LHR lift by 0.130 (95\% CI $[0.038, 0.247]$) and makes edits slightly more diffuse (concentration $-0.007$ and $-0.018$, fraction changed $+0.018$ and $+0.048$). The other differences, including CV and mean transfer in both directions, are within rerun noise; d0's p$\to$n CV drop of 0.140 (1 vs.\ 8 discordant images, $p=0.039$) has a CI that excludes zero but not the noise band. With 50 images and a single rerun pair as the noise reference, we read these as indications of sensitivity, not as a tuning result.

\begin{table*}[t]
\centering
\caption{Lock-width sensitivity, first 50 images per direction (LHR: 47 annotated). d2 is the committed setting, rerun. Differences are setting minus d2 rerun, with 95\% paired bootstrap CIs. Rerun noise is d2 rerun minus the committed run on the same images. $^{\dagger}$CI entirely outside $\pm$ rerun noise. $^{\ddagger}$CI excludes zero but not the rerun-noise band.}
\label{tab:suppl-lockwidth}
\footnotesize
\begin{tabular}{@{}lcccc@{}}
\toprule
Metric & d2 (rerun) & Rerun noise & d0 $-$ d2 & d4 $-$ d2 \\
\midrule
CV, n$\to$p & 0.960 & $0.000$ $[0.000, 0.000]$ & $-0.020$ $[-0.060, 0.000]$ & $0.000$ $[0.000, 0.000]$ \\
CV, p$\to$n & 0.740 & $-0.060$ $[-0.140, 0.000]$ & $-0.140$ $[-0.260, -0.040]$$^{\ddagger}$ & $+0.040$ $[-0.040, +0.120]$ \\
LHR lift (p$\to$n) & 2.489 & $-0.034$ $[-0.114, +0.043]$ & $+0.013$ $[-0.096, +0.124]$ & $-0.130$ $[-0.247, -0.038]$$^{\dagger}$ \\
Edit conc., n$\to$p & 0.248 & $0.000$ $[-0.002, +0.001]$ & $+0.028$ $[+0.025, +0.032]$$^{\dagger}$ & $-0.007$ $[-0.010, -0.004]$$^{\dagger}$ \\
Edit conc., p$\to$n & 0.317 & $-0.005$ $[-0.013, +0.002]$ & $+0.035$ $[+0.025, +0.043]$$^{\dagger}$ & $-0.018$ $[-0.025, -0.010]$$^{\dagger}$ \\
Frac.\ changed, n$\to$p & 0.664 & $0.000$ $[0.000, +0.001]$ & $-0.060$ $[-0.072, -0.050]$$^{\dagger}$ & $+0.018$ $[+0.008, +0.025]$$^{\dagger}$ \\
Frac.\ changed, p$\to$n & 0.629 & $+0.001$ $[-0.008, +0.009]$ & $-0.065$ $[-0.074, -0.056]$$^{\dagger}$ & $+0.048$ $[+0.037, +0.060]$$^{\dagger}$ \\
Transfer mean, n$\to$p & 0.262 & $-0.005$ $[-0.015, +0.002]$ & $-0.016$ $[-0.032, -0.001]$$^{\ddagger}$ & $+0.004$ $[-0.007, +0.015]$ \\
Transfer mean, p$\to$n & 0.251 & $-0.038$ $[-0.062, -0.014]$ & $-0.025$ $[-0.052, +0.003]$ & $+0.046$ $[+0.015, +0.077]$$^{\ddagger}$ \\
\texttt{nih}, n$\to$p & 0.042 & $-0.012$ $[-0.040, +0.010]$ & $-0.059$ $[-0.101, -0.018]$$^{\dagger}$ & $+0.017$ $[-0.010, +0.047]$ \\
\texttt{pc}, n$\to$p & 0.442 & $-0.002$ $[-0.007, +0.004]$ & $+0.019$ $[+0.006, +0.031]$$^{\dagger}$ & $-0.002$ $[-0.018, +0.012]$ \\
\texttt{mimic\_ch}, n$\to$p & 0.304 & $-0.002$ $[-0.006, +0.002]$ & $-0.008$ $[-0.024, +0.007]$ & $-0.004$ $[-0.013, +0.006]$ \\
\texttt{nih}, p$\to$n & 0.124 & $-0.060$ $[-0.105, -0.021]$ & $-0.038$ $[-0.082, +0.007]$ & $+0.073$ $[+0.018, +0.129]$$^{\ddagger}$ \\
\texttt{pc}, p$\to$n & 0.320 & $-0.004$ $[-0.040, +0.030]$ & $-0.029$ $[-0.070, +0.010]$ & $+0.033$ $[+0.001, +0.068]$$^{\ddagger}$ \\
\texttt{mimic\_ch}, p$\to$n & 0.307 & $-0.049$ $[-0.097, -0.006]$ & $-0.009$ $[-0.061, +0.045]$ & $+0.031$ $[-0.008, +0.070]$ \\
\bottomrule
\end{tabular}
\end{table*}
